\documentclass[10pt,a4paper]{amsart}
\usepackage{amsmath,amssymb,amsthm,graphicx}
% Retained preamble for the Stacks Project chapter snapshot in this directory.
% The chapter's own preamble.tex is a copy of the shared Stacks Project preamble; the block
% below is sliced MECHANICALLY, line for line, from it. Nothing was added, removed or
% redefined; the wrapper class, the xr-hyper/hyperref wiring and the externaldocument list
% are omitted because the assembler supplies the document class.
%
% lines taken from preamble.tex: 3, 5, 9, 11, 12, 156, 157, 158, ... (51 lines total)
%
% package coverage restored from the chapter's own preamble (2026-10-13): xy, xy-2cell

\usepackage[all]{xy}\usepackage{verbatim}
\xyoption{2cell}
\UseAllTwocells

\newenvironment{reference}{\comment}{\endcomment}
\newenvironment{slogan}{\comment}{\endcomment}
\newenvironment{history}{\comment}{\endcomment}
\theoremstyle{plain}
\newtheorem{theorem}[subsection]{Theorem}
\newtheorem{proposition}[subsection]{Proposition}
\newtheorem{lemma}[subsection]{Lemma}
\theoremstyle{definition}
\newtheorem{definition}[subsection]{Definition}
\newtheorem{example}[subsection]{Example}
\newtheorem{exercise}[subsection]{Exercise}
\newtheorem{situation}[subsection]{Situation}
\theoremstyle{remark}
\newtheorem{remark}[subsection]{Remark}
\newtheorem{remarks}[subsection]{Remarks}
\numberwithin{equation}{subsection}
\def\lim{\mathop{\mathrm{lim}}\nolimits}
\def\colim{\mathop{\mathrm{colim}}\nolimits}
\def\Spec{\mathop{\mathrm{Spec}}}
\def\Hom{\mathop{\mathrm{Hom}}\nolimits}
\def\Ext{\mathop{\mathrm{Ext}}\nolimits}
\def\SheafHom{\mathop{\mathcal{H}\!\mathit{om}}\nolimits}
\def\SheafExt{\mathop{\mathcal{E}\!\mathit{xt}}\nolimits}
\def\Sch{\mathit{Sch}}
\def\Mor{\mathop{\mathrm{Mor}}\nolimits}
\def\Ob{\mathop{\mathrm{Ob}}\nolimits}
\def\Sh{\mathop{\mathit{Sh}}\nolimits}
\def\NL{\mathop{N\!L}\nolimits}
\def\CH{\mathop{\mathrm{CH}}\nolimits}
\def\proetale{{pro\text{-}\acute{e}tale}}
\def\etale{{\acute{e}tale}}
\def\QCoh{\mathit{QCoh}}
\def\Ker{\mathop{\mathrm{Ker}}}
\def\Im{\mathop{\mathrm{Im}}}
\def\Coker{\mathop{\mathrm{Coker}}}
\def\Coim{\mathop{\mathrm{Coim}}}
\DeclareMathSymbol{\boxtimes}{\mathbin}{AMSa}{"02}
\def\QCohstack{\mathcal{QC}\!\mathit{oh}}
\def\Cohstack{\mathcal{C}\!\mathit{oh}}
\def\Spacesstack{\mathcal{S}\!\mathit{paces}}
\def\Quotfunctor{\mathrm{Quot}}
\def\Hilbfunctor{\mathrm{Hilb}}
\def\Curvesstack{\mathcal{C}\!\mathit{urves}}
\def\Polarizedstack{\mathcal{P}\!\mathit{olarized}}
\def\Complexesstack{\mathcal{C}\!\mathit{omplexes}}
\def\Pic{\mathop{\mathrm{Pic}}\nolimits}
\def\Picardstack{\mathcal{P}\!\mathit{ic}}
\def\Picardfunctor{\mathrm{Pic}}
\def\Deformationcategory{\mathcal{D}\!\mathit{ef}}
% lane-s3 patch-sd-envs.tex - DO NOT REMOVE LINES 16-21 (the executable patch).
% Contract (Chief-confirmed 2026-10-14):
%  * Lines 16-18: reference/slogan/history are declared as \comment containers in
%    every retained-preamble of this lane; the patch empties them so navigation
%    pointer text inside statements typesets without "Environment undefined".
%  * Line 21: \cite is globally no-op'd because the corpus rule is "cards carry
%    no bibliography": there is no bibliography backend at assembly time, and
%    every patched card's statement+proof span was scanned zero-cite before
%    patching (2026-10-14, lane-s3) - the only \cite occurrences live inside
%    the emptied reference environments (e.g. \cite{EGA} pointers). Removing
%    line 21 breaks compilation of every card whose statement carries a
%    reference environment with a citation (co01 wave: Citation `EGA' undefined).
%  * The patch is baked into task.json at card-build time (preamble_span
%    end_line must be the last line of this file - the builder computes it
%    dynamically since 2026-10-14); cards without patchenv never load it.
\renewenvironment{reference}{}{}
\renewenvironment{slogan}{}{}
\renewenvironment{history}{}{}
% empty out citations inside navigation-pointer environments (no bibliography
% backend; spans themselves are cite-free - scanned 2026-10-14, lane-s3)
\renewcommand{\cite}[2][]{}
\usepackage{cleveref}
\crefname{originallemma}{original lemma}{original lemmas}
\Crefname{originallemma}{Original lemma}{Original lemmas}
\makeatletter
\newlabel{algebraic-definition-viewed-as}{{04X6}{source}{}{sourcepointer-1}{}}
\newlabel{algebraic-definition-viewed-as@cref}{{[originallemma][0][]04X6}{[1][0][]source}{}{}{}}
\makeatother
\crefname{originallemma}{original lemma}{original lemmas}
\Crefname{originallemma}{Original lemma}{Original lemmas}
\makeatletter
\newlabel{algebraic-proposition-algebraic-stack-no-automorphisms}{{04SZ}{source}{}{sourcepointer-2}{}}
\newlabel{algebraic-proposition-algebraic-stack-no-automorphisms@cref}{{[originallemma][0][]04SZ}{[1][0][]source}{}{}{}}
\makeatother
\crefname{originallemma}{original lemma}{original lemmas}
\Crefname{originallemma}{Original lemma}{Original lemmas}
\makeatletter
\newlabel{lemma-thickening-diagonals}{{lemma-thickening-diagonals}{source}{}{sourcepointer-3}{}}
\newlabel{lemma-thickening-diagonals@cref}{{[originallemma][0][]lemma-thickening-diagonals}{[1][0][]source}{}{}{}}
\makeatother
\crefname{originallemma}{original lemma}{original lemmas}
\Crefname{originallemma}{Original lemma}{Original lemmas}
\makeatletter
\newlabel{spaces-more-morphisms-lemma-thickening-scheme}{{05ZR}{source}{}{sourcepointer-4}{}}
\newlabel{spaces-more-morphisms-lemma-thickening-scheme@cref}{{[originallemma][0][]05ZR}{[1][0][]source}{}{}{}}
\makeatother
\crefname{originallemma}{original lemma}{original lemmas}
\Crefname{originallemma}{Original lemma}{Original lemmas}
\makeatletter
\newlabel{stacks-morphisms-definition-absolute-separated}{{050D}{source}{}{sourcepointer-5}{}}
\newlabel{stacks-morphisms-definition-absolute-separated@cref}{{[originallemma][0][]050D}{[1][0][]source}{}{}{}}
\makeatother
\crefname{originallemma}{original lemma}{original lemmas}
\Crefname{originallemma}{Original lemma}{Original lemmas}
\makeatletter
\newlabel{stacks-morphisms-definition-separated}{{04YW}{source}{}{sourcepointer-6}{}}
\newlabel{stacks-morphisms-definition-separated@cref}{{[originallemma][0][]04YW}{[1][0][]source}{}{}{}}
\makeatother
\crefname{originallemma}{original lemma}{original lemmas}
\Crefname{originallemma}{Original lemma}{Original lemmas}
\makeatletter
\newlabel{stacks-morphisms-lemma-hierarchy}{{0AHJ}{source}{}{sourcepointer-7}{}}
\newlabel{stacks-morphisms-lemma-hierarchy@cref}{{[originallemma][0][]0AHJ}{[1][0][]source}{}{}{}}
\makeatother
\title{Thickening invariance of stack-level properties (Conrad)}
\author{The Stacks Project Authors}
\date{stacks-project:stacks-more-morphisms.tex:local-snapshot-20261012; GNU FDL 1.2 or later}
\begin{document}
\maketitle
\noindent\textit{Editorial note.} Statement span is lines 506-524 of sources/stacks-stacks-more-morphisms/pinned-original-source.tex; proof span is lines 526-557, the complete author proof (30 body lines).\par\medskip
\noindent\textit{Editorial note.} Navigation stubs: 7 cited-but-not-included labels are supplied as non-load-bearing pointers with their pinned tags where the snapshot tag index has them.\par\medskip
\noindent\textit{Editorial note.} extra\_packages is empty on purpose: the retained preamble already loads xy with options and a second declaration clashes in the assembler wrapper.\par\medskip
\noindent\textit{Editorial note.} Text is retained verbatim, including the author's own typos and placeholders.\par\medskip
\noindent\textit{Editorial note.} Queue index 36 of queue\_stacks-stacks-more-morphisms.json (pb=1241, stmt 506-524, proof 526-557).\par\medskip
\noindent\textit{Editorial note.} stacks-stacks-more-morphisms wave smm01 (2026-10-14): fresh adjudication accept per lane-s3/smm\_judgments.md; pool per Chief 10:5x open-pool ruling (net 59 with-proof rows, bloodline lane-s3 by anchor audit, v2 true-key namespace stacks:stacks-more-morphisms.tex, row lemma-deform-property-fp-over-ft avoided on cross-chapter key collision); 60 with-proof rows judged 45 accept / 14 reject / 1 avoid; patch-sd-envs contract spans scanned.\par\medskip
\noindent\textit{Editorial note.} Statement carries a reference/slogan/history environment (patchenv -> patch-sd-envs.tex).\par\medskip
\noindent\textit{Allowed standard prerequisites (editorial).}\begin{itemize}
\item Explicit prerequisite boundary. The cited inputs are prerequisites; the argument in the retained span is not replaced by any cited result.
\end{itemize}\medskip
\ifdefined\hypertarget\hypertarget{sourcepointer-1}{}\fi
\noindent\textit{Source reference.} Original lemma 04X6: sources/stacks-stacks-more-morphisms/pinned-original-source.tex, source lines 557--557. Cited by the retained statement or proof but not part of the collected spans; supplied only as a navigation pointer. The argument does not depend on its content.\par\medskip
\ifdefined\hypertarget\hypertarget{sourcepointer-2}{}\fi
\noindent\textit{Source reference.} Original lemma 04SZ: sources/stacks-stacks-more-morphisms/pinned-original-source.tex, source lines 557--557. Cited by the retained statement or proof but not part of the collected spans; supplied only as a navigation pointer. The argument does not depend on its content.\par\medskip
\ifdefined\hypertarget\hypertarget{sourcepointer-3}{}\fi
\noindent\textit{Source reference.} Original lemma lemma-thickening-diagonals: sources/stacks-stacks-more-morphisms/pinned-original-source.tex, source lines 454--468. Cited by the retained statement or proof but not part of the collected spans; supplied only as a navigation pointer. The argument does not depend on its content.\par\medskip
\ifdefined\hypertarget\hypertarget{sourcepointer-4}{}\fi
\noindent\textit{Source reference.} Original lemma 05ZR: sources/stacks-stacks-more-morphisms/pinned-original-source.tex, source lines 557--557. Cited by the retained statement or proof but not part of the collected spans; supplied only as a navigation pointer. The argument does not depend on its content.\par\medskip
\ifdefined\hypertarget\hypertarget{sourcepointer-5}{}\fi
\noindent\textit{Source reference.} Original lemma 050D: sources/stacks-stacks-more-morphisms/pinned-original-source.tex, source lines 557--557. Cited by the retained statement or proof but not part of the collected spans; supplied only as a navigation pointer. The argument does not depend on its content.\par\medskip
\ifdefined\hypertarget\hypertarget{sourcepointer-6}{}\fi
\noindent\textit{Source reference.} Original lemma 04YW: sources/stacks-stacks-more-morphisms/pinned-original-source.tex, source lines 557--557. Cited by the retained statement or proof but not part of the collected spans; supplied only as a navigation pointer. The argument does not depend on its content.\par\medskip
\ifdefined\hypertarget\hypertarget{sourcepointer-7}{}\fi
\noindent\textit{Source reference.} Original lemma 0AHJ: sources/stacks-stacks-more-morphisms/pinned-original-source.tex, source lines 557--557. Cited by the retained statement or proof but not part of the collected spans; supplied only as a navigation pointer. The argument does not depend on its content.\par\medskip
\begin{lemma}
\label{lemma-thickening-properties}
\begin{reference}
\cite[Theorem 2.2.5]{Conrad-moduli}
\end{reference}
Let $\mathcal{X} \subset \mathcal{X}'$ be a thickening of algebraic
stacks. Then
\begin{enumerate}
\item $\mathcal{X}$ is an algebraic space if and only if $\mathcal{X}'$
is an algebraic space,
\item $\mathcal{X}$ is a scheme if and only if $\mathcal{X}'$ is a scheme,
\item $\mathcal{X}$ is DM if and only if $\mathcal{X}'$ is DM,
\item $\mathcal{X}$ is quasi-DM if and only if $\mathcal{X}'$ is quasi-DM,
\item $\mathcal{X}$ is separated if and only if $\mathcal{X}'$ is separated,
\item $\mathcal{X}$ is quasi-separated if and only if $\mathcal{X}'$ is
quasi-separated, and
\item add more here.
\end{enumerate}
\end{lemma}

\begin{proof}
In each case we reduce to a question about the diagonal and then
we use Lemma \ref{lemma-thickening-diagonals} applied to the
morphism of thickenings
$$
(\mathcal{X} \subset \mathcal{X}') \to
\left(\Spec(\mathbf{Z}) \subset \Spec(\mathbf{Z})\right)
$$
We do this after viewing
$\mathcal{X} \subset \mathcal{X}'$ as a thickening of algebraic stacks
over $\Spec(\mathbf{Z})$ via
Algebraic Stacks, Definition \ref{algebraic-definition-viewed-as}.

\medskip\noindent
Case (1). An algebraic stack is an algebraic space if and only if its
diagonal is a monomorphism, see
Morphisms of Stacks, Lemma \ref{stacks-morphisms-lemma-hierarchy}
(this also follows immediately from Algebraic Stacks,
Proposition \ref{algebraic-proposition-algebraic-stack-no-automorphisms}).

\medskip\noindent
Case (2). By (1) we may assume that $\mathcal{X}$ and $\mathcal{X}'$
are algebraic spaces and then we can use
More on Morphisms of Spaces, Lemma
\ref{spaces-more-morphisms-lemma-thickening-scheme}.

\medskip\noindent
Case (3) -- (6). Each of these cases corresponds to a condition
on the diagonal, see Morphisms of Stacks, Definitions
\ref{stacks-morphisms-definition-separated} and
\ref{stacks-morphisms-definition-absolute-separated}.
\end{proof}

\end{document}
